%% ===================================================================
\section{Leading-log content of every term}\label{app:LLrows}
%% ===================================================================

In this appendix we list the LL limits of the amplitudes of Sec.~\ref{sec:xsec}, obtained from the factorized wave function~\eqref{eq:OmegaLLfact}. Squaring them as in Eqs.~\eqref{eq:CS1}--\eqref{eq:CS3} gives the LL content of every row.

\paragraph*{Conventions.} The measured gluon sits at $w$ in the amplitude and at $w'$ in its conjugate; its valence source is at $x$ ($x'$). The second gluon sits at $z$. In the real terms its source is at $y$ ($y'$); in the virtual terms it is emitted and absorbed by valence charges at $u$ and $v$. The colour $a$ of the measured gluon is in the final frame. The colour $h$ of the second gluon is in the frame before the target, like the index $b$ in $U^{ac}(w)\mf B_2^{bc}$. Repeated indices are summed, and all transverse positions other than $w,w'$ are integrated. The Wilson lines are general orthogonal matrices; the index order is significant. The LO amplitude is
\begin{equation}
\mathcal L^a_i(w)\equiv\KK^i(x-w)\big[U^{ab}(w)-U^{ab}(x)\big]\rho^b(x),\qquad {\rm LO}\propto\sum_{a,i}\big[\mathcal L^a_i(w')\big]^\dagger\mathcal L^a_i(w).
\end{equation}
$[X(w')]^\dagger$ means: take $X$, rename $w\to w'$, $x\to x'$, $y\to y'$, reverse the order of all charges and complex conjugate the numbers. Every row is normalized as
\begin{equation}
\frac{d^3N}{d^3k}\bigg|_{\rm row}=\frac{1}{(2\pi)^3}\,\frac{g^4}{8\pi^5k^+}\,\delta Y_{T\text{ or }P}\int_{w,w'}e^{-ik\cdot(w'-w)}\int[\text{Row}],
\label{eq:LLnorm}
\end{equation}
i.e.\ the LO prefactor $\frac{1}{(2\pi)^3}\frac{g^2}{2\pi^2k^+}$ times $\frac{\alpha_s}{\pi^2}$ per unit of rapidity. The rows are built from the pieces below as
\begin{align}
\text{Group I }(+{\rm c.c.}):&\quad \text{Row }X=\sum_{a,i}\big[\mathcal L^a_i(w')\big]^\dagger\,\mathfrak V^a_{X,i}(w)+{\rm c.c.},\notag\\
&\quad X\in\{{\rm I(a),I(b),II(a),II(b),III,IV,V,I}'\},\notag\\
\text{Group II}:&\quad \text{I}=[\mathfrak B]^\dagger\mathfrak B,\quad \text{II}=[\mathfrak X_1]^\dagger\mathfrak X_1,\quad \text{III}=[\mathfrak X_2]^\dagger\mathfrak X_2,\quad \text{IV}=[\mathfrak C]^\dagger\mathfrak C,\notag\\
\text{Group III }(+{\rm c.c.}):&\quad \text{I}=[\mathfrak X_1]^\dagger\mathfrak B,\ \ \text{II}=[\mathfrak X_2]^\dagger\mathfrak B,\ \ \text{III}=[\mathfrak C]^\dagger\mathfrak B,\notag\\
&\quad\text{IV}=[\mathfrak X_2]^\dagger\mathfrak X_1,\ \ \text{V}=[\mathfrak C]^\dagger\mathfrak X_1,\ \ \text{VI}=[\mathfrak C]^\dagger\mathfrak X_2,
\end{align}
where in Groups II and III $[\mathfrak Y]^\dagger\mathfrak Z\equiv\sum_{a,h,i,j}\int_z[\mathfrak Y^{ah}_{ij}(w',z)]^\dagger\mathfrak Z^{ah}_{ij}(w,z)$. In Group I, I(a) and I(b) are the normalization terms $-\bar{\mf A}\mf N_2$ and $U\bar{\mf N}_2\mf A$, II(a) is $U\mf A^{(3)}$, and II(b) is the complete $\bar{\mf A}^{(3)\dagger}_1$ (including the commutator with $\bar{\mf N}_2$). Rows III, IV and V then contain only the remaining term of Eqs.~\eqref{eq:VIII}--\eqref{eq:VV}, and I$'$ is the diagonal $\bar{\mf B}_3^\dagger$.

\subsection{Region T: \texorpdfstring{$\Lambda<p^+\ll k^+$}{Lambda<p+<<k+}}
\paragraph*{Two-gluon amplitudes.}
{\small
\begin{align}
\mathfrak B^{ah}_{ij}(w,z)=\;&\KK^{i}(x{-}w)\KK^{j}(y{-}z)\Big[U^{ab}(w)\rho^{h}(y)\rho^{b}(x)-U^{bh}(z)U^{bc}(y)U^{ad}(x)\rho^{c}(y)\rho^{d}(x)\Big]\notag\\
&+\KK^{i}(x{-}w)\KK^{j}(w{-}z)\Big[-if^{hbc}U^{ab}(w)\rho^{c}(x)+if^{bac}U^{bh}(z)U^{cd}(x)\rho^{d}(x)\Big]\label{eq:LLT2TB}
\\[4pt]
{\mathfrak X_1}^{ah}_{ij}(w,z)=\;&\KK^{i}(x{-}w)\KK^{j}(y{-}z)\Big[U^{bh}(z)U^{bc}(y)U^{ad}(x)\rho^{c}(y)\rho^{d}(x)-U^{bh}(z)U^{bc}(y)U^{ad}(w)\rho^{c}(y)\rho^{d}(x)\Big]\label{eq:LLT2TX1}
\\[4pt]
{\mathfrak X_2}^{ah}_{ij}(w,z)=\;&\KK^{i}(x{-}w)\KK^{j}(y{-}z)\Big[U^{ab}(x)U^{ch}(z)U^{cd}(y)\rho^{b}(x)\rho^{d}(y)-U^{ab}(x)\rho^{b}(x)\rho^{h}(y)\Big]\label{eq:LLT2TX2}
\\[4pt]
\mathfrak C^{ah}_{ij}(w,z)=\;&\KK^{i}(x{-}w)\KK^{j}(w{-}z)\Big[if^{bac}U^{bh}(z)U^{cd}(w)\rho^{d}(x)-if^{bac}U^{bh}(z)U^{cd}(x)\rho^{d}(x)\Big]\label{eq:LLT2TC}
\end{align}

}
\paragraph*{One-gluon (virtual) corrections.}
{\small
\begin{align}
{\mathfrak V_{\rm I(a)}}^{a}_{i}(w)=\;&\KK^{i}(x{-}w)\,K(u,v,z)\Big[\tfrac{1}{2}U^{ab}(x)\rho^{b}(x)\rho^{h}(u)\rho^{h}(v)\Big]\label{eq:LLT1TI_Nafter}
\\[4pt]
{\mathfrak V_{\rm I(b)}}^{a}_{i}(w)=\;&\KK^{i}(x{-}w)\,K(u,v,z)\Big[-\tfrac{1}{2}U^{bc}(u)U^{bd}(v)U^{ae}(w)\rho^{c}(u)\rho^{d}(v)\rho^{e}(x)\Big]\label{eq:LLT1TI_Nbarbefore}
\\[4pt]
{\mathfrak V_{\rm II(a)}}^{a}_{i}(w)=\;&\KK^{i}(x{-}w)\,K(u,v,z)\Big[-\tfrac{1}{2}U^{ab}(w)\rho^{h}(u)\rho^{h}(v)\rho^{b}(x)\Big]\notag\\
&+\KK^{i}(x{-}w)\,K(u,w,z)\Big[if^{hbc}U^{ab}(w)\rho^{h}(u)\rho^{c}(x)\Big]\notag\\
&+\KK^{i}(x{-}w)\,K(w,w,z)\Big[-\tfrac{3}{2}U^{ab}(w)\rho^{b}(x)\Big]\label{eq:LLT1TII_A3}
\\[4pt]
{\mathfrak V_{\rm II(b)}}^{a}_{i}(w)=\;&\KK^{i}(x{-}w)\,K(u,v,z)\Big[\tfrac{1}{2}U^{ab}(x)U^{cd}(u)U^{ce}(v)\rho^{b}(x)\rho^{d}(u)\rho^{e}(v)\Big]\label{eq:LLT1TII_A3bar}
\\[4pt]
{\mathfrak V_{\rm III}}^{a}_{i}(w)=\;&\KK^{i}(x{-}w)\,K(u,v,z)\Big[-U^{ab}(x)U^{ch}(z)U^{cd}(u)\rho^{b}(x)\rho^{d}(u)\rho^{h}(v)\Big]\label{eq:LLT1TIII}
\\[4pt]
{\mathfrak V_{\rm IV}}^{a}_{i}(w)=\;&\KK^{i}(x{-}w)\,K(u,v,z)\Big[U^{bh}(z)U^{bc}(u)U^{ad}(w)\rho^{c}(u)\rho^{h}(v)\rho^{d}(x)\Big]\notag\\
&+\KK^{i}(x{-}w)\,K(u,w,z)\Big[-if^{hbc}U^{dh}(z)U^{de}(u)U^{ab}(w)\rho^{e}(u)\rho^{c}(x)\Big]\label{eq:LLT1TIV}
\\[4pt]
{\mathfrak V_{\rm V}}^{a}_{i}(w)=\;&\KK^{i}(x{-}w)\,K(w,v,z)\Big[-if^{bcd}U^{eh}(z)U^{eb}(w)U^{ac}(w)\rho^{h}(v)\rho^{d}(x)\Big]\notag\\
&+\KK^{i}(x{-}w)\,K(w,w,z)\Big[-f^{bcd}f^{hde}U^{ac}(w)U^{gh}(z)U^{gb}(w)\rho^{e}(x)\Big]\label{eq:LLT1TV}
\\[4pt]
{\mathfrak V_{\rm I'}}^{a}_{i}(w)=\;&\KK^{i}(x{-}w)\,K(u,w,z)\Big[if^{bcd}U^{eg}(u)U^{ac}(w)U^{eb}(w)\rho^{g}(u)\rho^{d}(x)\Big]\notag\\
&+\KK^{i}(x{-}w)\,K(w,w,z)\Big[-\tfrac{3}{2}U^{ab}(w)\rho^{b}(x)\Big]\label{eq:LLT1TIa}
\end{align}

}

\subsection{Region P: \texorpdfstring{$k^+\ll p^+<\vee$}{k+<<p+<vee}}
Now the second gluon is the harder one, and the measured gluon can be emitted by it: this gives the terms with $\KK^i(z-w)$.
\paragraph*{Two-gluon amplitudes.}
{\small
\begin{align}
\mathfrak B^{ah}_{ij}(w,z)=\;&\KK^{i}(x{-}w)\KK^{j}(y{-}z)\Big[U^{ab}(w)\rho^{b}(x)\rho^{h}(y)-U^{bh}(z)U^{ac}(x)U^{bd}(y)\rho^{c}(x)\rho^{d}(y)\Big]\notag\\
&+\KK^{i}(z{-}w)\KK^{j}(y{-}z)\Big[-if^{bhc}U^{ab}(w)\rho^{c}(y)+if^{abc}U^{bh}(z)U^{cd}(y)\rho^{d}(y)\Big]\label{eq:LLP2PB}
\\[4pt]
{\mathfrak X_1}^{ah}_{ij}(w,z)=\;&\KK^{i}(x{-}w)\KK^{j}(y{-}z)\Big[U^{bh}(z)U^{bc}(y)U^{ad}(x)\rho^{c}(y)\rho^{d}(x)-U^{bh}(z)U^{bc}(y)U^{ad}(w)\rho^{c}(y)\rho^{d}(x)\Big]\label{eq:LLP2PX1}
\\[4pt]
{\mathfrak X_2}^{ah}_{ij}(w,z)=\;&\KK^{i}(x{-}w)\KK^{j}(y{-}z)\Big[U^{ab}(x)U^{ch}(z)U^{cd}(y)\rho^{b}(x)\rho^{d}(y)-U^{ab}(x)\rho^{b}(x)\rho^{h}(y)\Big]\label{eq:LLP2PX2}
\\[4pt]
\mathfrak C^{ah}_{ij}(w,z)=\;&\KK^{i}(z{-}w)\KK^{j}(y{-}z)\Big[-if^{abc}U^{bh}(z)U^{cd}(y)\rho^{d}(y)+if^{abc}U^{bh}(z)U^{cd}(z)\rho^{d}(y)\Big]\label{eq:LLP2PC}
\end{align}

}
\paragraph*{One-gluon (virtual) corrections.}
{\small
\begin{align}
{\mathfrak V_{\rm I(a)}}^{a}_{i}(w)=\;&\KK^{i}(x{-}w)\,K(u,v,z)\Big[\tfrac{1}{2}U^{ab}(x)\rho^{b}(x)\rho^{h}(u)\rho^{h}(v)\Big]\label{eq:LLP1PI_Nafter}
\\[4pt]
{\mathfrak V_{\rm I(b)}}^{a}_{i}(w)=\;&\KK^{i}(x{-}w)\,K(u,v,z)\Big[-\tfrac{1}{2}U^{bc}(u)U^{bd}(v)U^{ae}(w)\rho^{c}(u)\rho^{d}(v)\rho^{e}(x)\Big]\label{eq:LLP1PI_Nbarbefore}
\\[4pt]
{\mathfrak V_{\rm II(a)}}^{a}_{i}(w)=\;&\KK^{i}(x{-}w)\,K(u,v,z)\Big[-\tfrac{1}{2}U^{ab}(w)\rho^{b}(x)\rho^{h}(u)\rho^{h}(v)\Big]\label{eq:LLP1PII_A3}
\\[4pt]
{\mathfrak V_{\rm II(b)}}^{a}_{i}(w)=\;&\KK^{i}(x{-}w)\,K(u,v,z)\Big[\tfrac{1}{2}U^{bc}(u)U^{bd}(v)U^{ae}(x)\rho^{c}(u)\rho^{d}(v)\rho^{e}(x)\Big]\label{eq:LLP1PII_A3bar}
\\[4pt]
{\mathfrak V_{\rm III}}^{a}_{i}(w)=\;&\KK^{i}(x{-}w)\,K(u,v,z)\Big[-U^{bh}(z)U^{bc}(u)U^{ad}(x)\rho^{c}(u)\rho^{d}(x)\rho^{h}(v)\Big]\notag\\
&+\KK^{i}(z{-}w)\,K(u,v,z)\Big[if^{bch}U^{dc}(z)U^{de}(u)U^{ab}(z)\rho^{e}(u)\rho^{h}(v)\Big]\label{eq:LLP1PIII}
\\[4pt]
{\mathfrak V_{\rm IV}}^{a}_{i}(w)=\;&\KK^{i}(x{-}w)\,K(u,v,z)\Big[U^{bh}(z)U^{bc}(u)U^{ad}(w)\rho^{c}(u)\rho^{d}(x)\rho^{h}(v)\Big]\notag\\
&+\KK^{i}(z{-}w)\,K(u,v,z)\Big[-if^{bhc}U^{dh}(z)U^{de}(u)U^{ab}(w)\rho^{e}(u)\rho^{c}(v)\Big]\label{eq:LLP1PIV}
\end{align}

}
Rows V and I$'$ of Group I do not exist in region P: a soft measured gluon cannot absorb a harder gluon.
